% Answers to Chapter 2, from lectures/L02/appendix/c_solutions.md; the self-test from
% lectures/L02/README.md.

\facitavsnitt{c2}{Kapitel 2}{Algebra}
\begin{facit}

\svar{1.1} $R_{\text{TOT}} = 1\,\Omega$

\svar{1.2} \sv{a} $460\,\text{W}$ \sv{b} $40\,\text{W}$

\svar{1.3} $U_{\text{ut}} = 4\,\text{V}$

\svar{2.1} \sv{a} $3I_1 + 10I_2$ \sv{b} $3R_1 + 2R_2$ \sv{c} $7I^2 + 3I - 1$
\sv{d} Båda uttrycken ger $33$.

\svar{2.2} \sv{a} $U_R = RI_1 + RI_2 + RI_3$; varje term har enheten $\Omega \cdot \text{A} =
\text{V}$. \sv{b} $10\,\text{V}$ \sv{c} $2 + 3 + 5 = 10\,\text{V}$, samma svar.

\svar{2.3} \sv{a} $3I(2 + 3R)$ \sv{b} $(P_1 + P_2)(P_1 - P_2)$ \sv{c} $(U - 5)^2$
\sv{d} $(2R + 1)(2R - 1)$

\svar{2.4} \sv{a} $P_{\text{TOT}} = (R_1 + R_2 + R_3)I^2$ \sv{b} $1{,}25\,\text{W}$
\sv{c} $P_1 = 0{,}30\,\text{W}$, $P_2 = 0{,}45\,\text{W}$, $P_3 = 0{,}50\,\text{W}$, summan
$1{,}25\,\text{W}$. Den faktoriserade formen krävde minst arbete: en multiplikation med $I^2$ i
stället för tre.

\svar{2.5} \sv{a} Bryt ut $\frac{1}{R}$; kvar blir $U_1^2 - U_2^2$, en differens av två kvadrater,
som faktoriseras med konjugatregeln. \sv{b} $22\,\text{W}$
\sv{c} $P_1 = 72\,\text{W}$ och $P_2 = 50\,\text{W}$; skillnaden är $22\,\text{W}$.

\svar{2.6} \sv{a} $P = RI_0^2 + 2RI_0\,\Delta I + R\,\Delta I^2$ \sv{b} $44{,}1\,\text{W}$
\sv{c} $\Delta P = 4{,}1\,\text{W}$ \sv{d} $0{,}1\,\text{W}$, cirka $2{,}4\,\%$ av $\Delta P$.

\svar{2.7} \sv{a} $RI^2 = RI \cdot I$, och $RI = U$ är spänningen över motståndet enligt Ohms lag,
så $P = UI$. \sv{b} $R \cdot \frac{U^2}{R^2} = \frac{U^2}{R}$
\sv{c} $I = 2\,\text{A}$; alla tre formerna ger $24\,\text{W}$.

\svar{3.1} \sv{a} Tre \sv{b} $7$ respektive $-4$ \sv{c} $9$
\sv{d} $7$ och $R^2$ (eller $7$, $R$ och $R$) \sv{e} $29$

\svar{3.2} \sv{a} $3I + 14$ \sv{b} $U + 7$ \sv{c} $4R_1 + 3$ \sv{d} $5$

\svar{3.3} \sv{a} $R^2 + 8R + 15$ \sv{b} $2I^2 + 7I - 4$ \sv{c} $U^2 - 4$
\sv{d} $9R^2 - 12R + 4$ \sv{e} $7I$

\svar{3.4} \sv{a} $I^2 + 8I + 16$ \sv{b} $R^2 - 12R + 36$ \sv{c} $4U^2 + 12U + 9$
\sv{d} $25 - 10I + I^2$ \sv{e} $(3 + 4)^2 = 49$ men $3^2 + 4^2 = 25$; skillnaden är den dubbla
produkten $2ab = 24$.

\svar{3.5} \sv{a} $9996$ \sv{b} $1575$ \sv{c} $3{,}99$ \sv{d} $200$ \sv{e} $80\,\text{W}$

\svar{3.6} \sv{a} $4(2R + 3)$ \sv{b} $RI(I + 1)$ \sv{c} $(U + 7)(U - 7)$ \sv{d} $(I + 6)^2$
\sv{e} $(3R + 4)(3R - 4)$ \sv{f} $2(U + 2)(U - 2)$

\svar{3.7} \sv{a} $2I$ \sv{b} $R - 3$ \sv{c} $U + 2$ \sv{d} $I - 5$ \sv{e} $2R + 1$

\svar{3.8} \sv{a} $12\,\text{V}$ \sv{b} $18\,\text{W}$ \sv{c} $18\,\text{W}$, samma svar.
\sv{d} Effekten fyrdubblas, till $72\,\text{W}$: $R(2I)^2 = 4RI^2$.

\svar{3.9} \sv{a} Bråken har samma nämnare och täljarna summeras till nämnaren $R_1 + R_2$, så
$U_1 + U_2 = U_{\text{in}}$.
\sv{b} I kvoten tar $U_{\text{in}}$ och nämnaren $R_1 + R_2$ ut varandra, och kvar blir
$\frac{R_1}{R_2}$. \sv{c} $U_1 = \frac{U_{\text{in}}}{2}$
\sv{d} $U_1 = 3\,\text{V}$, $U_2 = 12\,\text{V}$

\svar{3.10} \sv{a} Den dubbla produkten saknas: $R^2 + 4R + 4$.
\sv{b} Faktorn ska multipliceras med båda termerna: $3I - 6$.
\sv{c} Ett bråk får bara förkortas med faktorer: $\frac{R}{4} + 1$.
\sv{d} Minustecknet byter tecken på alla termer: $-U + 3$.
\sv{e} Även variablerna multipliceras: $6R^2$.
\sv{f} Exponenten gäller hela parentesen: $4I^2$.

\svar[Testa dig själv]{T} \sv{1} $(3R + 2)(3R - 2)$ \sv{2} $5I + 7$
\sv{3} $20\,\text{W}$

\end{facit}
